\documentclass[11pt]{article}

\usepackage[a4paper,margin=1in]{geometry}
\usepackage{amsmath,amssymb,booktabs}
\usepackage{slashed}
\usepackage{hyperref}

\newcommand{\Bc}{B_c}
\newcommand{\Bcst}{B_c^\ast}
\newcommand{\BcOne}{B_{c1}}
\newcommand{\eps}{\epsilon}
\newcommand{\Slash}[1]{\slashed{#1}}
\newcommand{\GeV}{\mathrm{GeV}}

\title{Symbolic Projection Notes for \texorpdfstring{\(\BcOne\to \Bc^{(\ast)}\gamma\)}{Bc1 to Bc gamma}}
\author{Working notes}
\date{\today}

\begin{document}
\maketitle

\section{Scope}

These notes collect the compact FeynCalc projection results for the heavy-heavy
channels
\[
\BcOne(6743)\to \Bc\gamma,\qquad
\BcOne(6743)\to \Bcst\gamma,\qquad
\BcOne(6750)\to \Bc\gamma,\qquad
\BcOne(6750)\to \Bcst\gamma .
\]
The calculation is structurally different from the \(D_s\) and \(B_s\)
light-cone sum rules.  In the \(B_c\) system both valence constituents are
heavy, so the leading radiative OPE is a hard-photon heavy-heavy calculation.
There is no light spectator line to be replaced by photon distribution
amplitudes, and no \(\chi_s\langle\bar s s\rangle\) type term.

\section{Currents and Propagators}

For \(B_c^+\sim c\bar b\), we use
\[
j_5=\bar b\,i\gamma_5 c,\qquad
j_\rho^V=\bar b\,\gamma_\rho c ,
\]
and the two axial basis currents
\[
J_\mu^A=\bar b\,\gamma_\mu\gamma_5 c,\qquad
J_\mu^B=\frac{i}{m_b+m_c}\,
\bar b\,\sigma_{\mu\alpha}P^\alpha\gamma_5 c,
\qquad P=p+q .
\]
The physical currents are
\[
J_{1\mu}=\sin\theta\,J_\mu^A+\cos\theta\,J_\mu^B,\qquad
J_{2\mu}=\cos\theta\,J_\mu^A-\sin\theta\,J_\mu^B .
\]
This mixing is used at the invariant-amplitude level.  The OPE is first
derived for the basis currents because the two Dirac traces are independent;
the physical amplitudes are then
\[
\widehat\Pi_1=\sin\theta\,\widehat\Pi_A+\cos\theta\,\widehat\Pi_B,\qquad
\widehat\Pi_2=\cos\theta\,\widehat\Pi_A-\sin\theta\,\widehat\Pi_B .
\]
For the pseudoscalar final state,
\[
g_i=
\frac{m_b+m_c}{m_{\Bc}^2 f_{\Bc}\,m_i f_i}
\exp\!\left(\frac{m_i^2}{M_1^2}+\frac{m_{\Bc}^2}{M_2^2}\right)
\widehat\Pi_i^{(P)},
\]
while for the vector final state,
\[
g_i^\ast=
\frac{1}{m_{\Bcst}f_{\Bcst}\,m_i f_i}
\exp\!\left(\frac{m_i^2}{M_1^2}+\frac{m_{\Bcst}^2}{M_2^2}\right)
\widehat\Pi_i^{(V)}.
\]
The pilot numerical script applies this rotation.  The temporary approximation
is instead in the normalization: a common benchmark
\(f_{\BcOne}=0.373~\GeV\) is used for both physical axial-vector residues,
pending a separate \(f_1,f_2\) two-point extraction.

The free heavy-quark propagator in momentum space is
\[
S_Q^{(0)}(k)=\frac{i(\Slash{k}+m_Q)}{k^2-m_Q^2+i0},
\qquad Q=c,b .
\]
The hard photon is included by inserting the electromagnetic vertex on either
heavy line:
\[
S_Q^{(\gamma)}(k,k+q)=
S_Q^{(0)}(k+q)\,i e_Q\gamma_\nu\epsilon^\nu\,S_Q^{(0)}(k).
\]
For the \(\bar b\) line it is convenient to use the effective antiquark charge
\(e_{\bar b}=+e/3\), while \(e_c=+2e/3\).  Equivalently, one may keep
\(e_b=-e/3\) and track the orientation of the antiquark line.  The scripts use
\(e_{\bar b}\) explicitly.

Thus, yes: this is the same heavy-heavy propagator logic as in a \(B_c\)-mixing
sum-rule calculation.  The important difference from the strange-sector
radiative calculation is that we do not insert a light-quark photon DA.  If we
later include power corrections, the natural heavy-heavy correction is the
background-gluon expansion of the heavy propagators, e.g. gluon-condensate
terms, not strange-quark condensate photon DAs.

The submitted \(B_c\)-mixing work uses the same local basis,
\[
J_\mu^A=\bar b\gamma_\mu\gamma_5c,\qquad
J_\mu^B=\frac{i}{m_b+m_c}\bar b\sigma_{\mu\alpha}p^\alpha\gamma_5c,
\]
with \(m_b=4.18~\GeV\), \(m_c=1.27~\GeV\), and the final working window
\[
7~\GeV^2\le M^2\le 9~\GeV^2,\qquad
53~\GeV^2\le s_0\le 55~\GeV^2 .
\]
Those are useful starting values for the heavy-heavy OPE in the radiative
problem.  They are not yet automatically the final radiative-transition
windows; the final windows must be fixed by the three-point sum-rule stability,
pole dominance, and convergence tests.

\section{Step-by-Step Sum-Rule Construction}

This section is written as a checklist for reproducing the calculation.

\paragraph{1. Start from basis currents.}
The physical currents \(J_1,J_2\) are mixed combinations of \(J_A,J_B\).  We
derive the OPE for \(J_A\) and \(J_B\) separately only because the Dirac traces
are simpler.  The physical amplitudes are obtained afterward through the
\(\sin\theta,\cos\theta\) rotation.

\paragraph{2. Write the correlators.}
For \(B_c\gamma\),
\[
\Pi_\mu^{(5,K)}(p,q)=
i\int d^4x\,e^{ipx}
\langle\gamma(q,\epsilon)|
T\{\bar b(x)i\gamma_5c(x)\,J_{K\mu}^\dagger(0)\}|0\rangle ,
\qquad K=A,B .
\]
For \(B_c^\ast\gamma\),
\[
\Pi_{\mu\nu}^{(V,K)}(p,q)=
i\int d^4x\,e^{ipx}
\langle\gamma(q,\epsilon)|
T\{\bar b(x)\gamma_\nu c(x)\,J_{K\mu}^\dagger(0)\}|0\rangle .
\]

\paragraph{3. Contract the heavy fields.}
There are two hard-photon diagrams: photon emission from the \(c\) line and
photon emission from the \(\bar b\) line.  The free heavy propagator is
\[
S_Q^{(0)}(k)=\frac{i(\Slash{k}+m_Q)}{k^2-m_Q^2+i0},
\]
and the photon insertion is
\[
S_Q^{(\gamma)}(k,k+q)=
S_Q^{(0)}(k+q)\,i e_Q\Slash{\epsilon}\,S_Q^{(0)}(k).
\]
The charge convention in the scripts is \(e_c=+2/3\) and
\(e_{\bar b}=+1/3\).  No photon DA or light-quark condensate term appears in
this heavy-heavy leading baseline.

\paragraph{4. Evaluate and project the traces.}
For \(B_c\gamma\) project onto
\[
S_{\mu\nu}=p_\nu q_\mu-(p\cdot q)g_{\mu\nu},
\qquad
{\cal P}^{\mu\nu}=\frac{S^{\mu\nu}}{2(p\cdot q)^2},
\qquad
S_{\mu\nu}{\cal P}^{\mu\nu}=1 .
\]
For \(B_c^\ast\gamma\) the current diagnostic tensor is
\[
V_{\mu\nu\rho}=p_\nu\epsilon_{\mu\rho p q}
-(p\cdot q)\epsilon_{\mu\nu\rho p},
\qquad
V_{\mu\nu\rho}V^{\mu\nu\rho}=-4p^2(p\cdot q)^2 .
\]
The vector-channel result still needs matching to the physical
\(1^+\to1^-\gamma\) amplitude basis.

\paragraph{5. Reduce to denominator structures.}
For photon emission from the charm line use
\[
D_1=(k+q)^2-m_c^2,\quad D_2=k^2-m_c^2,\quad
D_3=(k-p)^2-m_b^2.
\]
For photon emission from the anti-bottom line use
\[
D_1=k^2-m_c^2,\quad D_2=(k-p)^2-m_b^2,\quad
D_3=(k-p+q)^2-m_b^2.
\]
The triangle pieces contain the full three-denominator product.  Terms where a
numerator cancels one denominator become two-denominator/contact terms and
must be included or shown to vanish.

\paragraph{6. Borel match and rotate.}
The true variables are \(p^2\) and \(P^2=(p+q)^2\), with
\[
p\cdot q=\frac{P^2-p^2}{2}.
\]
The pilot numerical calculation uses the diagonal choice \(M_1^2=M_2^2=2M^2\),
leading to the practical exponential
\[
\exp\left[\frac{m_i^2+m_f^2}{2M^2}\right].
\]
After obtaining \(\widehat\Pi_A,\widehat\Pi_B\), the physical combinations are
\[
\widehat\Pi_{6743}=\sin\theta\,\widehat\Pi_A+\cos\theta\,\widehat\Pi_B,
\qquad
\widehat\Pi_{6750}=\cos\theta\,\widehat\Pi_A-\sin\theta\,\widehat\Pi_B .
\]
For the \(B_c\gamma\) widths we use
\[
\Gamma(1^+\to0^-\gamma)=
\frac{\alpha_{\rm em}}{3}g^2E_\gamma^3,\qquad
E_\gamma=\frac{m_i^2-m_f^2}{2m_i}.
\]

\section{\texorpdfstring{\(\BcOne\to\Bc\gamma\)}{Bc1 to Bc gamma}: E1 Projection}

For the pseudoscalar final state we project onto
\[
S_{\mu\nu}=p_\nu q_\mu-(p\cdot q)g_{\mu\nu},
\qquad
{\cal P}_{\mu\nu}=\frac{S_{\mu\nu}}{2(p\cdot q)^2}.
\]
FeynCalc gives the normalization check
\[
S_{\mu\nu}{\cal P}^{\mu\nu}=1 .
\]
We denote
\[
p^2=p_2,\qquad p\cdot q=pq .
\]
After rewriting loop scalar products in terms of inverse denominators and
setting those denominators to zero, the projected triangle-core numerators are
\[
\begin{aligned}
C_A^{(c)} &=
-4i\,m_c+\frac{4i\,m_bm_c^2}{pq}-\frac{4i\,m_c^3}{pq},\\
C_A^{(\bar b)} &=
-4i\,m_b+\frac{4i\,m_b^3}{pq}-\frac{4i\,m_b^2m_c}{pq},
\end{aligned}
\]
so that
\[
C_A=e_c C_A^{(c)}+e_{\bar b}C_A^{(\bar b)} .
\]
For the tensor current,
\[
\begin{aligned}
C_B^{(c)} &=
\frac{4i\,m_bm_c-8i\,m_c^2}{m_b+m_c}
-\frac{4i\,m_c^2p_2}{(m_b+m_c)pq},\\
C_B^{(\bar b)} &=
-\frac{4i\,m_bm_c}{m_b+m_c}
-\frac{4i\,m_b^2p_2}{(m_b+m_c)pq},
\end{aligned}
\]
and
\[
C_B=e_c C_B^{(c)}+e_{\bar b}C_B^{(\bar b)} .
\]

The denominator-cancellation terms are stored in the corresponding output file.
They should
not be discarded automatically; they correspond to reduced two-denominator or
contact contributions and must be treated in the dispersion/Borel analysis.

\section{\texorpdfstring{\(\BcOne\to\Bcst\gamma\)}{Bc1 to Bcstar gamma}: Levi-Civita Projection}

For the vector final state a useful gauge-invariant structure is
\[
V_{\mu\nu\rho}
=p_\nu\,\eps_{\mu\rho p q}
-(p\cdot q)\,\eps_{\mu\nu\rho p}.
\]
It is transverse in the photon index for \(q^2=0\).  The FeynCalc checks are
\[
V_{\mu\nu\rho}V^{\mu\nu\rho}=-4p_2(pq)^2,\qquad
V_{\mu\nu\rho}{\cal P}^{\mu\nu\rho}=1 .
\]

The projected \(J_A\) triangle cores are
\[
\begin{aligned}
C_{A,V}^{(c)} &=
-\frac{4i\,m_bm_c}{p_2}
+\frac{2i\,m_c^2}{p_2}
+\frac{4i\,m_c^2}{pq},\\
C_{A,V}^{(\bar b)} &=
-\frac{2i\,m_b^2}{p_2}
+\frac{4i\,m_bm_c}{p_2}
+\frac{4i\,m_b^2}{pq}.
\end{aligned}
\]
The tensor-current vector cores are
\[
\begin{aligned}
C_{B,V}^{(c)} &=
\frac{2i\,m_c}{m_b+m_c}
+\frac{2i\,m_b^2m_c-4i\,m_bm_c^2+2i\,m_c^3}
       {(m_b+m_c)p_2} \\
&\quad
-\frac{4i\,m_bm_c^2-4i\,m_c^3}{(m_b+m_c)pq}
+\frac{2i\,m_c\,pq}{(m_b+m_c)p_2},\\
C_{B,V}^{(\bar b)} &=
\frac{2i\,m_b}{m_b+m_c}
+\frac{-6i\,m_b^3+4i\,m_b^2m_c+2i\,m_bm_c^2}
       {(m_b+m_c)p_2} \\
&\quad
-\frac{4i\,m_b^3-4i\,m_b^2m_c}{(m_b+m_c)pq}
+\frac{2i\,m_b\,pq}{(m_b+m_c)p_2}.
\end{aligned}
\]
Again,
\[
C_{A,V}=e_c C_{A,V}^{(c)}+e_{\bar b}C_{A,V}^{(\bar b)},\qquad
C_{B,V}=e_c C_{B,V}^{(c)}+e_{\bar b}C_{B,V}^{(\bar b)} .
\]

\section{First Numerical Hard-Photon Pass}

The script
\path{Bc_gamma/scripts/bc_radiative_pilot_sumrule.py}
uses the compact triangle cores above to produce a first hard-photon numerical
estimate.  For this pass we use the same central heavy-quark inputs and Borel
window as the submitted \(B_c\)-mixing analysis,
\[
m_b=4.18~\GeV,\qquad m_c=1.27~\GeV,\qquad
\theta=43.3^\circ,
\]
more precisely \(\theta_{\rm mean}=43.295^\circ\) with
\(\sigma_\theta=0.151^\circ\) and \(16\)--\(84\%\) interval
\([43.147^\circ,43.455^\circ]\) from the \(B_c\)-mixing Monte Carlo.  The
pilot table uses the rounded central value \(43.3^\circ\); the \(\theta\)
uncertainty is not yet included in the quoted radiative intervals.
\[
M^2=\{7,8,9\}~\GeV^2,\qquad s_0=\{53,54,55\}~\GeV^2.
\]
The decay-constant inputs are
\[
f_{\Bc}=0.371~\GeV,\qquad
f_{\Bcst}=0.384~\GeV,\qquad
f_{\BcOne}=0.373~\GeV .
\]
The first two are standard external benchmark inputs.  The last is a common
axial-vector benchmark used only for this first pass; separate physical
\(f_1\) and \(f_2\) residues have not yet been extracted from a two-point
mixed-current \(B_c\) sum rule in this folder.

\begin{table}[h]
\centering
\caption{Pilot hard-photon results in the \(B_c\) sector.  The intervals show
the envelope over \(M^2=\{7,8,9\}~\GeV^2\) and
\(s_0=\{53,54,55\}~\GeV^2\).}
\scriptsize
\setlength{\tabcolsep}{3pt}
\begin{tabular}{p{0.24\textwidth}p{0.18\textwidth}p{0.29\textwidth}p{0.15\textwidth}}
\toprule
Channel & Status & Coupling/projection coefficient & \(\Gamma\) (keV)\\
\midrule
\(\BcOne(6743)\to\Bc\gamma\)
& hard-QCDSR baseline
& \(+0.0219\,[0.0209,0.0232]~\GeV^{-1}\)
& \(0.108\,[0.098,0.121]\)\\
\(\BcOne(6750)\to\Bc\gamma\)
& hard-QCDSR baseline
& \(-0.354\,[-0.387,-0.331]~\GeV^{-1}\)
& \(29.6\,[25.9,35.3]\)\\
\(\BcOne(6743)\to\Bcst\gamma\)
& diagonal vector pilot
& \(-0.0620\,[-0.0669,-0.0583]\)
& \(46.3\,[41.0,53.8]\)\\
\(\BcOne(6750)\to\Bcst\gamma\)
& diagonal vector pilot
& \(+0.362\,[0.335,0.401]\)
& \(1.67\,[1.43,2.04]\times10^3\)\\
\bottomrule
\end{tabular}
\end{table}

The \(\Bc\gamma\) entries are the controlled hard-photon baseline from the
present calculation.  The \(\Bcst\gamma\) entries should be read as diagnostic
numbers only.  They are obtained from the gauge-invariant
\(V_{\mu\nu\rho}\) projection and a diagonal single-variable reduction.  Before
using them as final paper results one must match this coefficient to the
physical \(1^+\to1^-\gamma\) tensor basis and carry out the full double-Borel
analysis, including denominator-cancellation or contact terms.

\section{Comparison with Experiment and Other Theory}

LHCb has observed a wide peaking structure in the \(\Bc^+\gamma\) mass
spectrum with significance above seven standard deviations.  A minimal
two-peak description gives~\cite{LHCbBc1PObservation2025,LHCbBc1PStudy2025}
\[
M_1=6704.8\pm5.5\pm2.8\pm0.3~{\rm MeV},\qquad
M_2=6752.4\pm9.5\pm3.1\pm0.3~{\rm MeV},
\]
where the errors are statistical, systematic, and due to the \(B_c^+\) mass.
No individual radiative partial width is measured.  The natural widths are
expected to be only a few hundred keV and are neglected in the experimental
fit because they are small compared with the photon-energy resolution.

\begin{table}[h]
\centering
\caption{Qualitative comparison with the current experimental information
from LHCb~\cite{LHCbBc1PObservation2025,LHCbBc1PStudy2025}.}
\scriptsize
\setlength{\tabcolsep}{4pt}
\begin{tabular}{lll}
\toprule
Quantity & Experiment & Present input/result\\
\midrule
Observed \(B_c(1P)^+\)-like signal
& yes, \(>7\sigma\) in \(\Bc^+\gamma\)
& channels studied explicitly\\
Lower visible peak
& \(6704.8\pm5.5\pm2.8\pm0.3~{\rm MeV}\)
& input state \(6743~{\rm MeV}\)\\
Higher visible peak
& \(6752.4\pm9.5\pm3.1\pm0.3~{\rm MeV}\)
& input state \(6750~{\rm MeV}\)\\
Partial widths
& not measured
& \(\Bc\gamma\): \(0.108\), \(29.6~{\rm keV}\)\\
\bottomrule
\end{tabular}
\end{table}

The lower visible peak should not be identified one-to-one with the input
\(6743~{\rm MeV}\) state without care.  In the measured \(\Bc\gamma\) spectrum,
decays through \(\Bcst\gamma\), followed by an unreconstructed soft
\(\Bcst\to\Bc\gamma\), can shift visible peaks by the unknown hyperfine
splitting.

For comparison with other theory, Martín-González et al.~\cite{MartinGonzalez2022}
quote representative
quark-model E1 widths for the row labelled \(B_{c1}(1P)\):
\[
\Gamma[B_{c1}(1P)\to B_c\gamma]
  =1.5\times10^{-4}~{\rm keV},
\qquad
\Gamma[B_{c1}(1P)\to B_c^\ast\gamma]=146~{\rm keV}.
\]
Their table also lists earlier model values from
Refs.~\cite{Godfrey2004,Ebert2003,EichtenQuigg1994},
\[
\Gamma[B_{c1}(1P)\to B_c\gamma]=13,\ 18.4,\ 0.0~{\rm keV},
\]
\[
\Gamma[B_{c1}(1P)\to B_c^\ast\gamma]=60,\ 78.9,\ 99.5~{\rm keV}.
\]
The comparison therefore reads as follows.

\begin{table}[h]
\centering
\caption{Comparison with representative E1 model widths.  The model range is
collected from Ref.~\cite{MartinGonzalez2022}, which includes comparisons with
Refs.~\cite{Godfrey2004,Ebert2003,EichtenQuigg1994}.}
\scriptsize
\setlength{\tabcolsep}{4pt}
\begin{tabular}{p{0.23\textwidth}p{0.22\textwidth}p{0.25\textwidth}p{0.22\textwidth}}
\toprule
Channel & Other-theory scale & This work & Interpretation\\
\midrule
\(B_{c1}\to B_c\gamma\)
& \(0\) to \(18.4~{\rm keV}\)
& \(0.108~{\rm keV}\) for \(6743\), \(29.6~{\rm keV}\) for \(6750\)
& same broad scale, but state assignment and mixing differ\\
\(B_{c1}\to B_c^\ast\gamma\)
& \(60\) to \(146~{\rm keV}\)
& \(46.3~{\rm keV}\) for \(6743\), \(1.67~{\rm MeV}\) for \(6750\)
& lower pilot plausible; upper pilot flags unfinished vector normalization\\
\bottomrule
\end{tabular}
\end{table}

Thus the \(\Bc\gamma\) baseline is already useful for comparison, while the
\(\Bcst\gamma\) channel must still be completed with the physical tensor-basis
normalization and the full double-Borel/contact-term treatment.

\section{Next Step}

The double-dispersion variables satisfy
\[
P^2=p^2+2p\cdot q,\qquad pq=\frac{P^2-p^2}{2}.
\]
For \(\Bc\gamma\), the structure is close to the strange-sector E1 case, but
with two heavy masses.  For \(\Bcst\gamma\), the additional \(1/p_2\) and
\(pq/p_2\) structures mean that the double dispersion should be written
explicitly before numerical Borel analysis.  A diagonal single-variable
replacement can be useful as a calibration check, but it should not be the final
derivation for the vector final state.

\subsection*{Double-dispersion denominator}

For the photon emitted from the charm line, the three denominators are
\[
D_1=(k+q)^2-m_c^2,\qquad
D_2=k^2-m_c^2,\qquad
D_3=(k-p)^2-m_b^2 .
\]
With Feynman parameters \(x,y,z\), \(x+y+z=1\), the combined denominator is
\[
xD_1+yD_2+zD_3
=\ell^2-\Delta_c,
\]
where
\[
\Delta_c=(1-z)m_c^2+z m_b^2-zy\,p^2-xz\,P^2 .
\]
For photon emission from the \(\bar b\) line, using the momentum routing in the
FeynCalc scripts,
\[
D_1=k^2-m_c^2,\qquad
D_2=(k-p)^2-m_b^2,\qquad
D_3=(k-p+q)^2-m_b^2 ,
\]
and the analogous reduction gives, with \(x+y+z=1\),
\[
\Delta_{\bar b}
=x m_c^2+(1-x)m_b^2-x(1-x)p^2+xz(P^2-p^2).
\]
Equivalently one may relabel the Feynman parameters after choosing a different
loop momentum.  These forms make clear why the true radiative problem is a
double-dispersion problem in \(p^2\) and \(P^2\), not simply the two-point mixing
correlator with an extra charge factor.

\begin{thebibliography}{9}

\bibitem{LHCbBc1PObservation2025}
R. Aaij et al. (LHCb Collaboration),
``Observation of Orbitally Excited \(B_c^+\) States,''
Phys. Rev. Lett. 135, 231902 (2025).

\bibitem{LHCbBc1PStudy2025}
R. Aaij et al. (LHCb Collaboration),
``Study of \(B_c(1P)^+\) states in the \(B_c^+\gamma\) mass spectrum,''
Phys. Rev. D 112, 112003 (2025).

\bibitem{MartinGonzalez2022}
B. Martín-González, P. G. Ortega, D. R. Entem,
F. Fernández, and J. Segovia,
``Towards the discovery of novel \(B_c\) states: radiative and hadronic
transitions,''
arXiv:2205.05950.

\bibitem{Godfrey2004}
S. Godfrey,
Phys. Rev. D 70, 054017 (2004).

\bibitem{Ebert2003}
D. Ebert, R. N. Faustov, and V. O. Galkin,
Phys. Rev. D 67, 014027 (2003).

\bibitem{EichtenQuigg1994}
E. J. Eichten and C. Quigg,
Phys. Rev. D 49, 5845 (1994).

\end{thebibliography}

\end{document}
